\documentclass[10pt,a4paper]{article}
\usepackage[utf8]{inputenc}
\usepackage[T1]{fontenc}
\usepackage[english]{babel}
\usepackage{helvet}
\renewcommand{\familydefault}{\sfdefault}
\usepackage{amsmath,amssymb}
\usepackage{geometry}
\usepackage[table]{xcolor}
\usepackage{array,tabularx,booktabs}
\usepackage{hyperref}
\usepackage{fancyhdr}
\usepackage{microtype}
\usepackage{ragged2e}

\geometry{left=18mm,right=18mm,top=17mm,bottom=14mm,headheight=0pt,headsep=0pt,footskip=10mm}
\definecolor{navy}{HTML}{142A46}
\definecolor{teal}{HTML}{008891}
\definecolor{inkgray}{HTML}{536579}
\definecolor{lightgray}{HTML}{F1F5F8}
\definecolor{lightteal}{HTML}{E6F3F2}
\definecolor{linkblue}{HTML}{0875B0}
\hypersetup{colorlinks=true,urlcolor=linkblue,linkcolor=linkblue,pdfauthor={Chrystian Melo},pdftitle={CalcAreas and the ASPHAA Method - Technical Note},pdfsubject={Analytical comparison of constant-area and authalic methods}}
\setlength{\parindent}{0pt}
\setlength{\parskip}{3.0pt}
\setlength{\tabcolsep}{8pt}
\renewcommand{\arraystretch}{1.52}
\AtBeginDocument{\fontsize{8.7}{11.2}\selectfont\color{inkgray}}
\pagestyle{fancy}
\fancyhf{}
\renewcommand{\headrulewidth}{0pt}
\renewcommand{\footrulewidth}{0.3pt}
\fancyfoot[L]{\fontsize{7.3}{9}\selectfont\color{inkgray}Technical note supporting cv / interview}
\fancyfoot[R]{\fontsize{7.3}{9}\selectfont\bfseries\color{navy}1 / 1}

\newcommand{\sectionline}[2]{%
  \vspace{5pt}\noindent{\fontsize{10.8}{13}\selectfont\bfseries\color{navy}#1 \quad\textbar\quad #2}\par\vspace{1.8pt}%
}
\newcommand{\source}[1]{{\color{navy}\textbf{#1}}}
\newcommand{\callout}[2]{%
  \begingroup\setlength{\fboxsep}{9pt}%
  \noindent\colorbox{#1}{\parbox{\dimexpr\linewidth-2\fboxsep\relax}{\color{navy}#2}}\par
  \endgroup
}

\begin{document}
\thispagestyle{fancy}

{\color{teal}\rule{3pt}{10pt}\hspace{7pt}{\fontsize{9.5}{12}\selectfont\bfseries TECHNICAL NOTE / GEOSPATIAL COMPUTING}}\par
\vspace{1pt}
{\fontsize{19.5}{23}\selectfont\bfseries\color{navy}CalcAreas and the ASPHAA Method}\par
\vspace{3pt}
Mathematical foundations and reproducible comparison of areas in geographic grids (latitude/longitude).
\vspace{4pt}

\callout{lightgray}{Reference result at \textbf{50$^\circ$ latitude}: a \textbf{1$^\circ\times1^\circ$} cell estimated using the constant-area approximation is \textbf{55.21\% larger} than the corresponding WGS84 area. The authalic method removes this geometric bias in the analytical model.}

\sectionline{1}{Primary sources and parameters}
\source{CalcAreas [1]:} with \texttt{Use Authalic Area = false}, the documentation specifies a fixed distance of \textbf{111,255 m per degree} for both latitude and longitude, assigning the same area to every cell in the geographic grid. With \texttt{true}, it uses an authalic sphere.

\source{ASPHAA [2]:} Santini, Taramelli \& Sorichetta (2010), \textit{Transactions in GIS}, 14(3), 351--377. Area calculation using authalic latitude, preserving the area of the ellipsoid. \source{Erratum [3]:} corrections to equations (9)--(11), published in 2021.

\source{WGS84 reference ellipsoid [4]:} semi-major axis $a=6{,}378{,}137\ \mathrm{m}$; inverse flattening $1/f=298.257223563$; derived authalic radius $R_q\approx6{,}371{,}007.18\ \mathrm{m}$.

\sectionline{2}{Mathematical model}
Scenario: regular cells spanning $1^\circ$ in latitude $\times\ 1^\circ$ in longitude, centered at the indicated latitudes; arbitrary longitude. Boundaries follow meridians and parallels; angles in the formulas are expressed in radians. No land mask is applied.

\vspace{2pt}
\begingroup\setlength{\fboxsep}{8pt}
\noindent\colorbox{lightteal}{\parbox{\dimexpr\linewidth-2\fboxsep\relax}{%
\begin{minipage}[t]{0.485\linewidth}
{\color{navy}$\displaystyle A_{\mathrm{const}}=(111.255\,\Delta\varphi_{\mathrm{deg}})(111.255\,\Delta\lambda_{\mathrm{deg}})$}\par
For $1^\circ\times1^\circ$: $12{,}377.675\ \mathrm{km}^2$, regardless of latitude.
\end{minipage}\hfill
\begin{minipage}[t]{0.48\linewidth}
{\color{navy}$A_{\mathrm{ASPHAA}}=R_q^2\,\Delta\lambda\,\bigl(\sin\beta_N-\sin\beta_S\bigr)$}\par
$\sin\beta(\varphi)=q(\varphi)/q(90^\circ)$; $R_q^2=a^2q(90^\circ)/2$.
\end{minipage}%
}}\par
\endgroup

\vspace{2pt}
For geodetic latitude $\varphi$ and eccentricity $e^2=2f-f^2$:
\begingroup\color{navy}\begin{equation*}
q(\varphi)=(1-e^2)\left[\frac{\sin\varphi}{1-e^2\sin^2\varphi}-\frac{1}{2e}\ln\left(\frac{1-e\sin\varphi}{1+e\sin\varphi}\right)\right].
\end{equation*}\endgroup
\vspace{1pt}
Percentage error of the constant-area method: $100\times\bigl(A_{\mathrm{const}}/A_{\mathrm{WGS84}}-1\bigr)$. The WGS84 area was verified by numerical integration over the ellipsoidal surface; in the idealized model, it matches the authalic transformation.

\sectionline{3}{Calculated results -- $1^\circ\times1^\circ$ cells}
{\fontsize{8.25}{10}\selectfont
\rowcolors{2}{lightgray}{white}
\arrayrulecolor{white}
\begin{tabularx}{\linewidth}{>{\raggedright\arraybackslash}p{0.10\linewidth} >{\raggedright\arraybackslash}p{0.18\linewidth} >{\raggedright\arraybackslash}p{0.28\linewidth} >{\raggedright\arraybackslash}p{0.18\linewidth} X}
\rowcolor{navy}\color{white}\bfseries Latitude & \color{white}\bfseries Constant area & \color{white}\bfseries ASPHAA / WGS84 area & \color{white}\bfseries Constant-area error & \color{white}\bfseries ASPHAA error* \\
$0^\circ$  & 12,377.68 & 12,308.92 & +0.56\% & $\approx0\%$ \\
$30^\circ$ & 12,377.68 & 10,695.60 & +15.73\% & $\approx0\%$ \\
$45^\circ$ & 12,377.68 & 8,762.27 & +41.26\% & $\approx0\%$ \\
\rowcolor{lightteal}$50^\circ$ & 12,377.68 & 7,974.55 & +55.21\% & $\approx0\%$ \\
$60^\circ$ & 12,377.68 & 6,216.72 & +99.10\% & $\approx0\%$ \\
$70^\circ$ & 12,377.68 & 4,260.11 & +190.55\% & $\approx0\%$ \\
$80^\circ$ & 12,377.68 & 2,165.45 & +471.60\% & $\approx0\%$ \\
\end{tabularx}
}
\vspace{1pt}
{\fontsize{8}{10}\selectfont Area unit: km$^2$. *Errors of $\approx0\%$ follow from the mathematical equivalence between the authalic transformation and ellipsoidal area, not from a test run in Dinamica EGO.}

\sectionline{4}{Interpretation, metric limitations, and recommended wording}
Representative value: \textbf{+55.21\% at 50$^\circ$} (a difference of \textbf{4,403.13 km$^2$} per $1^\circ\times1^\circ$ cell). \textbf{Table maximum: +471.60\% at 80$^\circ$. Arithmetic mean across these 7 latitudes: +124.86\%;} this is an illustrative sample, not a global average or an actual distribution of cells.

Limitations: the error depends on latitude and resolution; 471.60\% is not a universal maximum. This calculation compares theoretical methods and does not measure execution time, empirical code accuracy, or average improvement in environmental simulations. The percentage should not be presented as a throughput gain or as the error corrected across all cells.

\vspace{2pt}
\callout{lightteal}{\textbf{Résumé bullet:} Implemented scientific algorithms for spatial analysis, including authalic-sphere grid-cell area computation based on published GIS methods, addressing latitude-dependent area distortions exceeding 50\% near 50$^\circ$ latitude.}

\sectionline{5}{References and traceability}
{\fontsize{8.3}{10.4}\selectfont
[1] CalcAreas -- official documentation: \href{https://www.csr.ufmg.br/dokuwiki/doku.php?id=calc_areas}{https://www.csr.ufmg.br/dokuwiki/doku.php?id=calc\_areas}\par
[2] ASPHAA -- paper (2010): \href{https://doi.org/10.1111/j.1467-9671.2010.01200.x}{https://doi.org/10.1111/j.1467-9671.2010.01200.x}\par
[3] Paper erratum (2021): \href{https://doi.org/10.1111/tgis.12785}{https://doi.org/10.1111/tgis.12785}\par
[4] WGS84 parameters -- NGA: \href{https://earth-info.nga.mil/index.php?action=wgs84\&dir=wgs84}{https://earth-info.nga.mil/index.php?action=wgs84\&dir=wgs84}
}
\end{document}
